\scanpage{70}
des autom des $(C,J)$ est donc \unsure{can$^{\mathrm{t}}$.|can$^{\mathrm{t}}$,}
\[
:\quad \unsure{\mathfrak S|\mathscr S}_\tau\times\mathfrak S_K\times\mathfrak S_{K'}\quad(\simeq\mathfrak S_2\times\mathfrak S_3\times\mathfrak S_3]
\]
d'ordre $72$.

\uline{NB} Il y a une involution naturelle dans cette catégorie des $(\tau,K,K')$, savoir
\[
(\tau,K,K')\longmapsto(\tau,K',K)
\]
Elle se traduit en associant à $(C,J)$ le $(C,J')$, où $J'$ est l'ens de cardinal $3$ des éléments de $C$ liés à tous les él. de $J$. Dans le cas où
\[
J=\{\unsure{\delta_{12}|J_{12}|\xi_{12}},\delta_{23},\delta_{31}\},\quad\text{on a}\quad J'=\{\delta_{45},\delta_{56},\delta_{\unsure{61|64}}\}.
\]

\uline{Corollaire.} $J$ appartient à \unsure{exactmt|exactement} $2$ bases, \deleted{(dans le cas envisagé, $\{\delta_{12},\delta_{23},\delta_{\unsure{34|31}}$ \unsure{prolongé|prolongés} soit} \inserted{savoir $J\cup(b\vert K')$ et $J\cup(b'\vert K')$} (dans \unsure{l'exp.|l'ex.} le \unsure{voisin|voisinage|réunion} de $J=\{\delta_{12},\delta_{23},\delta_{31}\}$ avec $\{\delta_4,\delta_5,\delta_6\}$ ou $\{\delta'_4,\delta'_5,\delta'_6\}$) dont l'intersection est réduite à $J$. \deleted{Ainsi}

Considérons alors l'application $\{\beta',\beta''\}\longrightarrow\beta'\cap\beta''$ de l'ens des couples de bases telles que $\mathrm{card}\,\beta'\cap\beta''=\underline{3}$, vers l'ens des parties $3$ libres de $C$, cette application est encore bijective. Cet ens s'identifie aussi à l'ens des paires $\{r,s\}$ de racines telles que $\widehat{r,s}=60^\circ$.
\margin{\unsure{expliciter|explicitez} position relative de $2$ racines}
% slanted
% archive sheet number: 69
